\section{Known Open Issues}
\label{sec:open-issues}

Snapshot as of 22 July 2026, the end of Session 5. This mirrors the ``Known
Open Issues'' section of \texttt{docs/SESSION\_CONTEXT.md}; kept here so this
report is self-contained.

\begin{enumerate}[label=\textbf{I\arabic*.}, leftmargin=1.6cm]

  \item \textbf{Session 5 geometry fixes are not yet hardware-verified
    end-to-end.} \texttt{\_fix\_object\_height} / \texttt{\_fix\_release\_height}
    (Section~\ref{sec:tech-decisions}) and the grasp-width edge-inset fix are
    verified against real captured hardware \texttt{debug/} data and
    unit-level replay, but not through a fresh, live end-to-end hardware run
    since landing. \textbf{Action:} re-run the \texttt{sort\_hard} and
    \texttt{arith\_hard} prompts that originally surfaced these bugs and
    confirm no regressions -- listed first in the September restart plan
    (Executive Summary).

  \item \textbf{The nudge-outside logic only corrects a drifted release
    after the fact.} \texttt{\_fix\_release\_height} fixes the height and
    position of a release that drifts into a zone it was not meant for, but
    does nothing to stop the LLM from \emph{choosing} that drifted position
    in the first place. This remains open; see Roadmap item R7 (exposing
    geometry fixes as callable skills) as one possible direction.

  \item \textbf{\texttt{VLM\_REASONING\_EFFORT='none'} still not A/B tested
    on hardware.} Carried over from Session 4, unchanged. Somewhat
    superseded in relevance by the benchmark's finding
    (Section~\ref{sec:benchmark}) that the dominant VLM failure mode is
    arithmetic/reasoning under load rather than raw grounding precision,
    which this setting was aimed at -- worth re-evaluating whether it is
    still the right lever before spending time A/B testing it.

  \item \textbf{Plane-induced-homography $\to$ depth-based top-down
    rectification rewrite: not started.} Carried over unchanged from earlier
    sessions; no immediate plan attached.

  \item \textbf{Zone size is unknown in plain VLM (point-grounding) mode.}
    \texttt{distribute\_zone\_releases} spaces out colliding releases using
    only the released object's \texttt{grasp\_width}, with no notion of the
    target zone/tray's actual size -- nothing stops the spacing grid from
    growing past a tray's physical edge. See Roadmap item R1.

  \item \textbf{Groq empty-completion failures on \texttt{fhp}/\texttt{ffhp}.}
    (\texttt{qwen3.6-27b}) -- hypothesized to be the model exhausting its
    reasoning-token budget on the longer chained predicates$\to$plan prompt;
    \texttt{cot\_sc}'s shorter single-shot prompt does not show this. A
    reduced-schema / split-grounding-from-planning mitigation was discussed
    (Session 3) but not implemented.

  \item \textbf{Nebius bad-grounding-depth failures} on
    \texttt{fhp}/\texttt{ffhp}/\texttt{always\_act} -- the plan parses fine
    (not a JSON error) but returned pixel coordinates do not land on the
    target's actual surface, so deprojection gets no valid depth. Reproduced
    across multiple reasoning methods on Nebius; same schema-complexity
    hypothesis as the Groq issue above, but manifesting as a
    grounding-accuracy failure rather than a parse failure.

  \item \textbf{Physical-setup-dependent calibration will drift over the
    pause.} ChArUco board pose (\texttt{board\_in\_base\_*}) and workspace
    limits in \texttt{scene\_mock.json} are tied to the exact physical
    mounting of the camera and table; a month of disuse is enough time for a
    bump, a re-mount, or a lighting change to invalidate them silently.
    \textbf{Action:} re-verify calibration (Executive Summary, restart
    checklist item 2) before trusting any pick/place coordinates in the
    first September session.

\end{enumerate}
